\section{The recurrent logic unit}
\label{sec:unit}

\subsection{Variant A: one logic block}
Fig.~\ref{fig:logic-block} shows the unit. The LUT has four input wires $x_1,\dots,x_4$. In the hidden layer the MUX selects the register output, so the block output $y$ is the stored state, and $y$ is routed back to one LUT input. With $x_1$ the incoming spike bit and $x_2=y$ the unit's own previous output, the next state is
\begin{equation}
y_{t+1} = T\bigl[x_{1,t},\,y_t,\,x_{3,t},\,x_{4,t}\bigr],
\end{equation}
where $T$ is the truth table. With $x_3,x_4$ unused, $T$ has four rows, i.e.\ four learned bits and $2^4=16$ possible units. The row index is the change between old and new bit (Table~\ref{tab:rows}); the stored bit says what the unit does in response.

\begin{table}[t]
\centering
\caption{Rows of the two-input unit. $x_1$ is the new input bit, $x_2=y$ the previous output; the last column lists the four learned bits (rows $00,01,10,11$) for four example behaviors.}
\label{tab:rows}
\begin{tabular}{@{}cll@{\qquad}ll@{}}
\toprule
$x_1x_2$ & Meaning of the row & & Example unit & Stored bits\\
\midrule
00 & stays low & & follower ($y=x_1$) & 0\,0\,1\,1\\
01 & falling & & hold ($y=x_2$) & 0\,1\,0\,1\\
10 & rising & & latch high ($x_1\lor x_2$) & 0\,1\,1\,1\\
11 & stays high & & toggle ($x_1\oplus x_2$) & 0\,1\,1\,0\\
\bottomrule
\end{tabular}
\end{table}

\paragraph{MUX modes.} With the MUX on $Q$ the block is a clocked state machine, which is the hidden-layer use. With the MUX on the LUT output the block is stateless; such a block must never be fed back into itself, because that would form a combinational loop with no clock. The free inputs $x_3,x_4$ can carry two more upstream spikes or a leak tick (Section~\ref{sec:recover}).

\subsection{Variant B: logic cluster}
A cluster contains several LUTs that process subsets or combinations of signals, state logic that combines their outputs, registers that retain state, muxes that select or route signals, and recurrent feedback. A LUT is a component of the cluster, not necessarily the neuron itself. A general abstraction is
\begin{equation}
s_{t+1} = F\bigl(s_t,\, b_1(x_t),\dots,b_K(x_t),\, q_t\bigr),
\end{equation}
where $s_t$ is the internal state, $b_i$ are branch LUT functions and $q_t$ are selected or routed signals. The output may be a spike, a state-change signal, a transition code or a multi-bit state, depending on the experiment. This is a conceptual starting point, not a finalized neuron definition. Variant A is the case $K=1$ with a two-bit state. \todo{decide which variant is the primary unit of the paper}

\paragraph{State representation.} An illustrative two-bit encoding is 00 stable/low, 01 falling, 10 rising, 11 stable/high. State semantics and bit encoding are separate design variables; Gray, one-hot or learned encodings, and interpretations such as phase, event history or persistence, are all to be explored.

\paragraph{State-dependent receptive field.} The current state can drive muxes that choose which incoming signals or branch functions the cluster examines at the next update, so the unit does not evaluate a fixed input subset every cycle. This has to be compared with static selection under equal hardware budgets.

\subsection{Layers}
Fig.~\ref{fig:architecture} shows the proposed network.
\begin{description}[leftmargin=1.5em,style=nextline]
\item[Input layer (multi-bit).] A delta encoder keeps the previous sample, subtracts and compares with a threshold, emitting signed events on two wires (up, down) per input. Alternatives are thermometer coding (bit $i$ is $x>\tau_i$, so a level crossing becomes a rise or fall in the unit watching that level) and sigma-delta coding (a one-bit stream whose average is the value). A time replay counter reproduces the differential-time code of~\cite{windhager2024} in clock time (Section~\ref{sec:recover}).
\item[Hidden layers (logic blocks).] Rows of units as in Fig.~\ref{fig:logic-block}. Connectivity is set by place-and-route through the switch boxes, so wiring is chosen when the design is built. Each layer adds one clock of delay.
\item[Output layer (multi-bit).] Counters tally the spikes of each class unit over $T$ steps and the largest count wins; alternatively an up/down counter reconstructs a value from rise and fall events. Carry chains suit both.
\end{description}

\paragraph{Relation to the paper's layer model.} The baseline processes one layer after another over stored spike trains. Clock-locked recurrent units cannot be hosted by that scheme; Section~\ref{sec:ledger} (recurrence) discusses what to keep, drop or replace.

\begin{figure}[!htbp]
\centering
\resizebox{0.9\linewidth}{!}{%
\begin{tikzpicture}[font=\small,>=Stealth,line width=0.7pt]
% LUT
\node[draw,fill=black!8,minimum width=2.6cm,minimum height=3cm,align=center] (lut) at (0,0) {Look-up\\table\\(LUT)};
% inputs x1, x3, x4 from the layer below / neighbours
\foreach \yy/\lab in {1.1/x_1,-0.37/x_3,-1.1/x_4}{
  \node[anchor=east] at (-2.7,\yy) {$\lab$};
  \draw (-2.6,\yy) -- (-1.3,\yy);
}
% x2 is the feedback wire
\node at (-2.2,0.62) {$x_2$};
\draw (7.4,0) -- (7.4,-2.6) -- (-3.4,-2.6) -- (-3.4,0.37) -- (-1.3,0.37);
\node[font=\footnotesize,anchor=north] at (2.0,-2.6) {feedback of $y$ through the routing fabric};
% LUT output
\draw (1.3,0) -- (2.3,0);
\fill (2.3,0) circle (2pt);
% bypass to MUX
\draw (2.3,0) -- (2.3,0.6) -- (6.0,0.6);
% FF
\draw (2.3,0) -- (2.3,-0.6) -- (3.2,-0.6);
\node[draw,minimum width=1.6cm,minimum height=1.8cm] (ff) at (4.0,-0.6) {FF};
\node[anchor=west,font=\footnotesize] at (3.2,-0.3) {D};
\node[anchor=east,font=\footnotesize] at (4.8,-0.3) {Q};
\draw (4.8,-0.6) -- (6.0,-0.6);
% clock
\draw (2.9,-2.1) -- (2.9,-1.2) -- (3.2,-1.2);
\draw (3.2,-1.05) -- (3.4,-1.2) -- (3.2,-1.35);
\node[font=\footnotesize,anchor=north] at (2.9,-2.1) {clk};
% MUX
\draw (6.0,1.0) -- (6.0,-1.0) -- (6.8,-0.6) -- (6.8,0.6) -- cycle;
\node[rotate=90,font=\footnotesize] at (6.4,0) {MUX};
\draw (6.8,0) -- (7.4,0);
\fill (7.4,0) circle (2pt);
\draw[->] (7.4,0) -- (8.3,0) node[right] {$y$};
\end{tikzpicture}}
\caption{Logic block used as a recurrent unit (schematic redrawn after the generic LUT--FF--MUX logic cell). The MUX selects the registered path $Q$; the block output $y$ is routed back to LUT input $x_2$, so the LUT sees the new input bit $x_1$ together with the unit's own previous output.}
\label{fig:logic-block}
\end{figure}

\begin{figure}[!htbp]
\centering
\resizebox{0.95\linewidth}{!}{%
\begin{tikzpicture}[
  font=\small,
  >=Stealth,
  cont/.style={draw,dashed,rounded corners=4pt,black!60},
  tealbox/.style={draw=teal!60!black,fill=teal!12,rounded corners=3pt,align=center,minimum height=1.2cm},
  purpbox/.style={draw=violet!70!black,fill=violet!10,rounded corners=3pt,align=center,minimum height=1.2cm},
  greybox/.style={draw=black!50,fill=black!6,rounded corners=3pt,align=center,minimum height=1.2cm},
  arr/.style={-{Stealth[length=2mm]},thick}
]
% ---------------- input layer
\draw[cont] (0,0) rectangle (14,-2.6);
\node[anchor=west,font=\small\bfseries] at (0.3,-0.4) {Input layer: multi-bit logic};
\node[tealbox,minimum width=4.4cm] (enc) at (3.0,-1.6) {\textbf{Delta encoder}\\{\footnotesize up / down spike wires}};
\node[tealbox,minimum width=4.4cm] (rep) at (8.4,-1.6) {\textbf{Time replay counter}\\{\footnotesize restores spike timing}};
\draw[arr] (enc.east) -- (rep.west);
% ---------------- hidden layer
\draw[cont] (0,-3.4) rectangle (14,-10.6);
\node[anchor=west,font=\small\bfseries] at (0.3,-3.8) {Hidden layer: logic blocks};
\foreach \x/\n in {1.9/d1,5.3/d2,8.7/d3,12.1/d4}{
  \node[purpbox,minimum width=3.0cm] (\n) at (\x,-5.2) {\textbf{Dendrite}\\{\footnotesize stateless LUT4}};
}
\node[purpbox,minimum width=4.0cm] (mrg) at (7,-7.2) {\textbf{Merge LUT4}\\{\footnotesize combines 4 branches}};
\node[purpbox,minimum width=4.6cm,line width=0.9pt] (soma) at (7,-9.2) {\textbf{Soma block}\\{\footnotesize LUT + FF, MUX on Q}};
\node[greybox,minimum width=3.4cm] (tick) at (2.2,-9.2) {\textbf{Leak tick}\\{\footnotesize shared slow pulse}};
\draw[arr] (rep.south) -- (8.4,-4.6);
\draw[arr] (d1.south) -- (6.1,-6.6);
\draw[arr] (d2.south) -- (6.7,-6.6);
\draw[arr] (d3.south) -- (7.3,-6.6);
\draw[arr] (d4.south) -- (7.9,-6.6);
\draw[arr] (mrg.south) -- (soma.north);
\draw[arr] (tick.east) -- (soma.west);
\draw[arr] (9.3,-8.95) -- (10.3,-8.95) -- (10.3,-9.55) -- (9.3,-9.55);
\node[anchor=west,font=\footnotesize] at (10.5,-9.25) {$y\rightarrow x_2$};
% ---------------- output layer
\draw[cont] (0,-11.6) rectangle (14,-14.2);
\node[anchor=west,font=\small\bfseries] at (0.3,-12.0) {Output layer: multi-bit logic};
\node[tealbox,minimum width=4.0cm] (cnt) at (6.4,-13.2) {\textbf{Spike counter}\\{\footnotesize carry chain, per class}};
\node[tealbox,minimum width=4.0cm] (arg) at (11.2,-13.2) {\textbf{Argmax}\\{\footnotesize largest count wins}};
\draw[arr] (cnt.east) -- (arg.west);
\draw[arr] (soma.south) -- (7,-11.1) -- (6.4,-11.1) -- (cnt.north);
\end{tikzpicture}}
\caption{Block architecture. The input and output layers use multi-bit logic. The hidden layer uses only LUT/FF/MUX blocks: stateless dendrite and merge LUTs (MUX on the LUT output) feed a soma block (MUX on the registered output) with self-feedback and a shared leak tick. Lateral links between somas are not drawn.}
\label{fig:architecture}
\end{figure}


\begin{table}[t]
\centering
\caption{Baseline~\cite{windhager2024} versus the proposed hidden layer.}
\label{tab:compare}
\small
\begin{tabular}{@{}p{3.1cm}p{5.7cm}p{5.7cm}@{}}
\toprule
Aspect & Windhager et al. & This work\\
\midrule
Hidden unit & LIF: multi-bit potential, shifter decay, adder, comparator & LUT--FF(--MUX) state machine\\
Parameters & 4--6-bit weights rotated through one shared datapath per layer & truth-table bits and wiring\\
Next-spike-time stage & common stage per layer & none: all units update every clock\\
Time handling & differential time code, processing decoupled from spike time & clock-locked; decoupling restorable (Sec.~\ref{sec:recover})\\
Connectivity & fully connected, no links inside a layer & sparse fan-in; lateral and recurrent links possible\\
Readout & spike count, largest wins & same\\
Training & surrogate gradient (snnTorch) & open (Sec.~\ref{sec:learning})\\
\bottomrule
\end{tabular}
\end{table}

\FloatBarrier
